%% 
%% Template article for cas-sc documentclass for 
%% double column output.

\documentclass[a4paper,fleqn]{cas-sc}

% If the frontmatter runs over more than one page
% Use the longmktitle option.

%\documentclass[a4paper,fleqn,longmktitle]{cas-sc}

%\usepackage[numbers]{natbib}
%\usepackage[authoryear]{natbib}
\usepackage[numbers]{natbib}
%% The amssymb package provides various useful mathematical symbols
\usepackage{amssymb}
%% The amsmath package provides various useful equation environments.
\usepackage{amsmath}
\usepackage{epstopdf}
\usepackage{mathptmx,mathtools}
\usepackage{subcaption}
\usepackage{float} % Add this in your preamble
\usepackage[export]{adjustbox}
\usepackage{array} % Needed for m{} column type
\usepackage[percent]{overpic} % <-- put this in your preamble



\usepackage{amsmath}
\usepackage{nicefrac}
\usepackage{dirtytalk}
\usepackage{graphicx}   % For including images
\usepackage{tikz}
\usetikzlibrary{calc,tikzmark} % needed <<<<<<<<<<
\usepackage{xcolor}     % For coloring
\usepackage{lineno}
\renewcommand\linenumberfont{\normalfont\bfseries\small} % <-- Increase line number size
\usetikzlibrary{arrows.meta}
\usetikzlibrary{calc, positioning}
%%%Author macros
\def\tsc#1{\csdef{#1}{\textsc{\lowercase{#1}}\xspace}}
\tsc{WGM}
\tsc{QE}

\begin{document}
	
	\setcounter{figure}{1}
	
		%%%%%%%%%%%%%%%%%%%%%%%%%%%%%________Figure_4___________%%%%%%%%%%%%%%%%%%%%%
	\begin{figure}[!t]
		\centering
		
		% -------- Subfigure (a) -------- %
		\begin{subfigure}[t]{0.52\textwidth}
			\centering
			\begin{overpic}[width=\textwidth,trim={9cm 1cm 10cm 1cm},clip]{Fig4a.eps}
				\put(1,70){\fontfamily{ptm}\selectfont\textbf{(a)}}
				\put(47,70){\fontfamily{ptm}\selectfont\textbf{(b)}}
				\put(1,33){\fontfamily{ptm}\selectfont\textbf{(c)}}
				\put(47,33){\fontfamily{ptm}\selectfont\textbf{(d)}}
			\end{overpic}
		\end{subfigure}%
		% -------- Subfigure (b) -------- %
		\begin{subfigure}[t]{0.51\textwidth}
			\centering
			\begin{overpic}[width=\textwidth,trim={1cm 0cm 21cm 0cm},clip]{Fig4b.eps}
				\put(5,70){\fontfamily{ptm}\selectfont\textbf{(e)}}
				
				
				\color{black}
				\linethickness{0.5pt}
				\put(30.5,17.5){\dashbox{1}(45.1,0){}} % horizontal dashed box (acts as dashed line)
				\put(30.5,41.5){\dashbox{1}(45.1,0){}} % horizontal dashed box (acts as dashed line)
				\put(30.5,41.5){\color{red}\circle*{1.5}}
				\put(30.5,17.5){\color{red}\circle*{1.5}}
				
				\color{black}
				
				\put(60,28){\rotatebox{90}{\fontfamily{ptm}\fontsize{5}{5}\selectfont $L_w$}}
				
				
				\put(45,42){\fontfamily{ptm}\fontsize{5}{5}\selectfont{Front stagnation (bubble)}}
				
				
				\put(49,15){\fontfamily{ptm}\fontsize{5}{5}\selectfont{Rear stagnation (wake)}}
				\put(18,27){\fontfamily{ptm}\fontsize{5}{5}\selectfont{Circulating}}
				
				\color{red}
				\put(12,11){\fontfamily{ptm}\fontsize{5}{5}\selectfont{Outermost}}
				\color{blue}
				\put(32,19){\fontfamily{ptm}\fontsize{5}{5}\selectfont{Innermost}}
				
			\end{overpic}
		\end{subfigure}
		
		\caption{Primary wake of a rising bubble visualized using PRA fields for case $C_6$ at , highlighting regions of rotational coherence represented by elliptic LCSs, color-coded with the stretching factor $\lambda$. The solid line denotes the outermost elliptic LCS, while the dashed lines indicate the inner ones. Corresponding forward FTLE at  fields illustrating material transport barriers in the wake, shown as red solid lines. (c) Enlarged Length of the primary wake measured from the outermost hyperbolic LCS and compared with the experimental correlation of~\citet{Bhaga1980Exp}.}
		
		\label{Fig4}
	\end{figure}
	%%%%%%%%%%%%%%%%%%%%%%%%%%%%%________Figure_4_end___________%%%%%%%%%%%%%%%%%%%%%
	
\end{document}